\section{Quặng và quy trình luyện kim}

\subsection{Độ tinh khiết và hiệu suất}

Bài toán luyện kim thường có hai lớp dữ kiện trước khi tính theo PTHH:
\[
\boxed{m_{\text{quặng}}}
\rightarrow
\boxed{m_{\text{hợp chất có ích}}}
\rightarrow
\boxed{n}
\rightarrow
\boxed{\text{các PTHH}}
\rightarrow
\boxed{H}
\rightarrow
\boxed{m_{\text{kim loại}}}.
\]
Nếu mỗi công đoạn có hiệu suất riêng, phải áp dụng đúng hiệu suất tại công đoạn đó; không gộp tùy ý nếu đề không cho phép.

\subsection{Quặng sulfide và sản phẩm phụ}

Quặng sulfide thường được nung trong không khí để chuyển sulfide thành oxide, đồng thời tạo \ce{SO2}. Khí \ce{SO2} có thể được thu hồi để sản xuất sulfuric acid:
\[
\ce{2SO2 + O2 ->[V2O5,t^\circ] 2SO3},
\]
\[
\ce{SO3 + H2O -> H2SO4}.
\]
Như vậy một bài luyện kim có thể đồng thời hỏi kim loại chính và sản phẩm phụ.

\begin{vidu}
Một quặng chứa $80\%$ \ce{ZnS} theo khối lượng. Dùng \SI{1.50}{\ton} quặng, các phản ứng tạo Zn xảy ra hoàn toàn. Tính khối lượng Zn thu được. Cho $M_{\ce{ZnS}}=97$, $M_{\ce{Zn}}=65$ g/mol.
\end{vidu}
\begin{loigiai}
\textbf{Bước 1. Tính lượng chất có ích}
\[
m_{\ce{ZnS}}=1.50\times0.80=\SI{1.20}{\ton}.
\]

\textbf{Bước 2. Viết sơ đồ phản ứng}
\[
\ce{ZnS -> ZnO -> Zn}.
\]
Tỉ lệ mol tổng quát là $n_{\ce{ZnS}}:n_{\ce{Zn}}=1:1$.

\textbf{Bước 3. Tính giá trị trung gian}
\[
m_{\ce{Zn},\text{lt}}=1.20\times\frac{65}{97}\approx\SI{0.804}{\ton}.
\]

\textbf{Bước 4. Kết luận}
\textbf{Vậy,} khối lượng Zn thu được khoảng \SI{0.804}{\ton}.
\end{loigiai}

\subsection{Hai nguồn nguyên liệu -- lập hệ phương trình}

Nếu hai loại quặng cùng cung cấp một kim loại nhưng có thành phần khác nhau, có thể đặt khối lượng hai loại quặng là $x,y$. Một phương trình đến từ tổng khối lượng nguyên liệu; phương trình còn lại từ lượng kim loại hoặc sản phẩm thu được.

\begin{ghinho}{Nguyên tắc}
\[
\text{số ẩn} \leftrightarrow \text{số dữ kiện độc lập}.
\]
Nếu có hai ẩn, cần hai phương trình độc lập. Nếu có ba ẩn và đủ ba dữ kiện, cách làm vẫn tương tự; không xuất hiện phương pháp hóa học mới.
\end{ghinho}

\begin{tuluyen}
Hai loại quặng A, B chứa lần lượt $70\%$ và $50\%$ \ce{Fe2O3}. Cần phối trộn tổng cộng \SI{100}{\ton} quặng để lượng \ce{Fe2O3} là \SI{62}{\ton}. Tính khối lượng mỗi loại quặng.
\dapso{A: \SI{60}{\ton}; B: \SI{40}{\ton}.}
\begin{loigiaiBT}
Gọi khối lượng quặng A, B lần lượt là $x,y$ (tấn).
\[
\begin{cases}
x+y=100,\\
0.70x+0.50y=62.
\end{cases}
\Rightarrow
\begin{cases}
x=60,\\
y=40.
\end{cases}
\]

\textbf{Vậy,} cần \SI{60}{\ton} quặng A và \SI{40}{\ton} quặng B.
\end{loigiaiBT}
\end{tuluyen}
